\documentclass[a4paper,12pt]{article}
\usepackage[T2A]{fontenc}
\usepackage[utf8]{inputenc}
\usepackage[russian]{babel}
\usepackage{amsmath,amssymb,mathtools}
\usepackage{helvet}
\renewcommand{\familydefault}{\sfdefault}
\usepackage[margin=24mm,top=22mm,bottom=24mm]{geometry}
\usepackage{microtype}
\usepackage{xcolor}
\usepackage{tikz}
\usepackage{booktabs,tabularx,array}
\usepackage{enumitem}
\usepackage{fancyhdr}
\usepackage{hyperref}
\definecolor{accent}{HTML}{008080}
\definecolor{darkaccent}{HTML}{004D4D}
\definecolor{textgray}{HTML}{333333}
\definecolor{muted}{HTML}{647575}
\hypersetup{colorlinks=true,linkcolor=darkaccent,pdfauthor={Лёвин Артём Александрович},pdftitle={Ревью домашней работы — София Лузянина — 23.11.26}}
\setlength{\parindent}{0pt}
\setlength{\parskip}{7pt}
\setlength{\emergencystretch}{2em}
\setlist[itemize]{leftmargin=17pt,topsep=2pt,itemsep=2pt}
\pagestyle{fancy}
\fancyhf{}
\fancyhead[L]{\small\color{muted}Ревью домашней работы}
\fancyhead[R]{\small\color{muted}23.11.2026}
\fancyfoot[C]{\footnotesize\color{textgray}Лёвин Артём Александрович · эксклюзивно для Софии Лузяниной}
\fancyfoot[R]{\footnotesize\color{muted}\thepage}
\setlength{\headheight}{14pt}
\renewcommand{\headrulewidth}{0.25pt}
\renewcommand{\footrulewidth}{0pt}
\newcommand{\sectionline}{\par\vspace{2mm}\noindent\color{accent}\rule{\linewidth}{0.65pt}\color{black}\par\vspace{1mm}}
\newcommand{\lbl}[1]{\textbf{\color{darkaccent}#1}}
\begin{document}
\thispagestyle{empty}
\vspace*{24mm}
{\color{accent}\large ФИЗИКА\par}
\vspace{19mm}
{\Huge\bfseries Ревью домашней\\[3mm] работы\par}
\vspace{8mm}
{\Large\color{darkaccent}Кинематика. Движение и броски тел\par}
\vspace{32mm}
{\Large София Лузянина\par}
\vspace{5mm}
{\large 23 ноября 2026 г.\par}
\vfill
{\normalsize Лёвин Артём Александрович\par}
{\small\color{muted}Эксклюзивно для Софии Лузяниной\par}
\vspace{15mm}
\begin{tikzpicture}[x=1cm,y=1cm]
\draw[accent!55,line width=0.95pt] (0,0.4) .. controls (4,1.5) and (9,-0.5) .. (15.8,0.65);
\end{tikzpicture}
\newpage
\section*{Сводная таблица ошибок}
\sectionline
Проверены десять заданий по кинематике. В таблицу включено только задание, в котором обнаружена неточность в обосновании. Правильно решённые задания в ревью не приводятся.
\par\vspace{4mm}
{\small
\renewcommand{\arraystretch}{1.35}
\begin{tabularx}{\linewidth}{@{}p{13mm}X p{34mm}p{34mm}@{}}
\toprule
\textbf{№} & \textbf{Что исправить} & \textbf{Ваш результат} & \textbf{Правильный результат}\\
\midrule
\hyperref[task:9]{\textbf{9}} & Знак вертикального ускорения в уравнении высоты; обоснование времени полёта & Численные ответы верны & Те же значения; требуется исправить ход решения \\
\bottomrule
\end{tabularx}
}\par
\vspace{7mm}
\lbl{Уточнение об исходных данных.} В условии задачи ускорение свободного падения явно не указано. В Вашем решении использовано $g=10\,\text{м/с}^2$; ниже сохраняем это стандартное приближение.
\newpage
\thispagestyle{empty}
\null\vfill
\begin{center}
{\color{accent}\Large\bfseries Разбор ошибок}\par
\vspace{5mm}
{\color{muted}Точечная коррекция решения и самостоятельное закрепление}
\end{center}
\vfill\null
\newpage
\phantomsection\label{task:9}
\section*{Задание 9. Бросок под углом к горизонту}
\sectionline
\lbl{Условие.} Тело брошено с земли под углом $45^\circ$ к горизонту со скоростью $20\,\text{м/с}$ и упало на тот же уровень. Найдите время подъёма, полное время полёта, максимальную высоту и дальность. Примем $g=10\,\text{м/с}^2$.

\lbl{Ваш ответ.} $t_{\text{под}}=\sqrt2\,\text{с}$; $t_{\text{пол}}=2\sqrt2\,\text{с}$; $h_{\max}=10\,\text{м}$; $L=40\,\text{м}$. Все четыре значения правильные.

\lbl{Последний правильный шаг.} Вы разложили начальную скорость на горизонтальную и вертикальную составляющие и верно использовали формулу дальности полёта:
\[
 v_{0x}=v_0\cos\alpha,\qquad v_{0y}=v_0\sin\alpha,\qquad
 L=\frac{v_0^2\sin 2\alpha}{g}=40\,\text{м}.
\]

\textcolor{red!75!black}{\textbf{Ошибка в обосновании.}} В записи высоты появляется соотношение
\[
 h(t)=v_{0y}t+\frac{gt^2}{2}.
\]
Если вертикальная ось направлена вверх, а символ $g$ обозначает положительную величину $10\,\text{м/с}^2$, знак перед квадратом времени должен быть \textbf{отрицательным}. Далее в решении знак при преобразовании уравнения меняется без пояснения. В результате правильные ответы получены при внутренне противоречивой записи формулы.

\lbl{Почему это важно.} На движущееся тело действует ускорение свободного падения, направленное вниз. При положительном направлении оси $Oy$ вверх проекция ускорения равна $a_y=-g$. Поэтому уравнение координаты имеет вид
\[
 y(t)=y_0+v_{0y}t+\frac{a_yt^2}{2}
     =y_0+v_{0y}t-\frac{gt^2}{2}.
\]
Положительный знак перед $gt^2/2$ означал бы ускорение \emph{вверх}, что противоречит выбранной системе координат.

\lbl{Исправление.} Высота отсчитывается от точки броска, поэтому $y_0=0$ и
\[
 \boxed{h(t)=v_0\sin\alpha\,t-\frac{gt^2}{2}}.
\]
Вся дальнейшая работа должна опираться на эту формулу.

\newpage
\section*{Полное правильное решение задания 9}
\sectionline
\lbl{1. Проекции начальной скорости.} Поскольку $\sin45^\circ=\cos45^\circ=\sqrt2/2$, получаем
\[
 v_{0x}=20\cdot\frac{\sqrt2}{2}=10\sqrt2\,\text{м/с},
 \qquad v_{0y}=10\sqrt2\,\text{м/с}.
\]
\lbl{2. Время подъёма.} В верхней точке траектории вертикальная составляющая скорости равна нулю:
\[
 v_y(t)=v_{0y}-gt,\quad
 0=10\sqrt2-10t,\quad
 \boxed{t_{\text{под}}=\sqrt2\,\text{с}}.
\]
\lbl{3. Полное время полёта.} Возвращение на исходный уровень означает $h(t)=0$:
\[
 0=10\sqrt2\,t-5t^2
   =5t(2\sqrt2-t).
\]
Корень $t=0$ соответствует началу движения. Момент приземления:
\[
 \boxed{t_{\text{пол}}=2\sqrt2\,\text{с}}.
\]
\lbl{4. Максимальная высота.} Подставляем время подъёма в уравнение высоты:
\[
 h_{\max}=10\sqrt2\cdot\sqrt2-5(\sqrt2)^2
 =20-10=\boxed{10\,\text{м}}.
\]
\lbl{5. Дальность.} Горизонтальная скорость постоянна:
\[
 L=v_{0x}t_{\text{пол}}
 =10\sqrt2\cdot2\sqrt2
 =\boxed{40\,\text{м}}.
\]
\lbl{Проверка.} При $t=0$ и при $t=2\sqrt2\,\text{с}$ получаем $h=0$. В верхней точке $v_y=0$, а $h=10\,\text{м}$. Результаты согласованы с траекторией.

\textcolor{darkaccent}{\textbf{Итог.}} Все четыре найденных Вами значения верны. Исправить необходимо запись вертикального ускорения и уравнения высоты, чтобы решение стало физически последовательным.

\vspace{2mm}
\section*{Задача на закрепление}
Мяч брошен с земли под углом $30^\circ$ к горизонту с начальной скоростью $20\,\text{м/с}$. Примите $g=10\,\text{м/с}^2$. Направьте ось $Oy$ вверх. Запишите уравнение $y(t)$ и найдите время подъёма, полное время полёта, максимальную высоту и дальность. Покажите, какой знак имеет проекция ускорения на ось $Oy$.
\end{document}
